(1) 固定 \(m\) 时上自旋数为 \(N_+=N(1+m)/2\)，微观态数 \(\Omega=N!/(N_+!N_-!)\)。Stirling 公式给每自旋熵 \(s(m)=-k_B[p\ln p+(1-p)\ln(1-p)]\)，其中 \(p=(1+m)/2\)，端点采用 \(0\ln0=0\)。能量每自旋为 \(e(m)=-Jm^2/2-\mu Bm\)，所以 \(f=e-Ts=-Jm^2/2-\mu Bm+k_BT[p\ln p+(1-p)\ln(1-p)]\)。对 \(|m|<1\) 求导得到 \(f'(m)=-Jm-\mu B+k_BT\operatorname{artanh}m\)。平衡要求 \(f'=0\) 且选全局最小自由能，故 \(m=\tanh[(Jm+\mu B)/(k_BT)]\)，不是把外场直接代入无相互作用的 Brillouin 函数。

(2) 零场下 \(m=0\) 总是自洽解，但其稳定性由自由能二阶导数而非仅由方程存在性决定。计算 \(f''(m)=-J+k_BT/(1-m^2)\)，因此 \(f''(0)=k_BT-J\)，临界温度为 \(T_c=J/k_B\)。若 \(T>T_c\)，零磁化是唯一稳定的全局极小；若 \(T<T_c\)，零磁化成为局域极大，两侧出现一对由 \(m_0=\tanh[(T_c/T)m_0]\) 确定的稳定解 \(\pm m_0\)。在 \(T=T_c\) 二次曲率消失，必须保留四次项判断稳定性，不能据此宣称自由能完全平坦。

(3) 在 \(T>T_c\) 且 \(B\to0\) 时，利用 \(\tanh z=z+O(z^3)\) 得 \(m=(Jm+\mu B)/(k_BT)+O(B^3)\)，于是 \(m=\mu B/(k_BT-J)+O(B^3)\)，\(\chi_B=N\mu^2/[k_B(T-T_c)]\)。这是对磁场 \(B\) 的导数而非无量纲 SI 磁化率，量纲须保留。对 \(B=0\)、\(T\lesssim T_c\)，把自洽式展开为 \(m=(T_c/T)m-\tfrac13(T_c/T)^3m^3+O(m^5)\)，除去非零 \(m\) 后得到 \(m_0^2=3[(T_c/T)-1]/(T_c/T)^3+O((T_c-T)^2)\simeq3(T_c-T)/T_c\)，故自发磁化的临界指数为 \(1/2\)。

(4) 对小 \(m\) 展开熵项得到 \(f=-k_BT\ln2+\tfrac12(k_BT-J)m^2+\tfrac1{12}k_BT m^4-\mu Bm+O(m^6)\)。在 \(T=T_c\) 时驻值方程的首项是 \(Jm^3/3=\mu B\)，所以 \(m\sim\operatorname{sgn}(B)[3\mu|B|/J]^{1/3}\)，临界等温指数为 \(3\)。在零场、\(T<T_c\) 邻近临界点，极小化四阶多项式得 \(m^2\simeq3(J-k_BT)/(k_BT)\)，代回得到有序相相对于顺磁解析延拓的自由能差 \(\Delta f=-3(J-k_BT)^2/(4k_BT)\simeq-3k_B(T_c-T)^2/(4T_c)\)。因此 \(C=-NT\partial_T^2 f_{\rm eq}\) 的临界下侧比上侧多 \(3Nk_B/2\)。这是一种平均场热力学极限的非解析跃变；有限 \(N\) 的配分函数解析，不存在真正的奇点。
